% ===================================================================
%  Αρχείο: 18713.tex (Fragment)
%  Αυτόματη μετατροπή μέσω doc2latex.py
% ===================================================================

ΘΕΜΑ 4

Δίνεται το πολυώνυμο $P(x) = 2x^3 - \alpha x^2 + 2x + \beta$, $\alpha, \beta \in \mathbb{R}$. Αν $P(1) = 2$ και το υπόλοιπο της διαίρεσης $P(x) : (x - 2)$ ισούται με 15,

\begin{alist}
\item Να δείξετε ότι $P(x) = 2x^3 - x^2 + 2x - 1$.

\monades{8}

\item

\begin{rlist}
\item Να δείξετε ότι το πολυώνυμο $\pi(x) = x^2 + 1$ είναι παράγοντας του $P(x)$.

\monades{5}

\item Να λύσετε την εξίσωση $P(x) = 0$.

\monades{5}

\end{rlist}
\item Να λύσετε την εξίσωση $\sigma\upsilon\nu^3 x + \frac{1}{2}\sigma\upsilon\nu x = 1 - \eta\mu^2 x$, $x \in (0, 2\pi)$.

\monades{7}
\end{alist}
